% TOP_PROBLEM: {"problemId": "problem.greenberg-s-conjecture-on-vanishing-of-iwasawa-invariants-of-totally-real-fields", "canonicalTitle": "Greenberg's Conjecture on Vanishing of Iwasawa Invariants of Totally Real Fields", "releaseRank": 428, "primaryDomain": "number_theory_arithmetic_geometry", "primaryDomainLabel": "Number theory & arithmetic geometry", "displayStatus": "open_with_solved_subcases", "releaseStatus": "open_with_solved_subcases"}
% !TeX program = lualatex
% ATTEMPT_STATUS: unresolved
% Research continuation by GPT_6_Astra_Ultra, 27 September 2026.
\documentclass[11pt]{article}
\usepackage[margin=25mm]{geometry}
\usepackage{fontspec}
\setmainfont{Latin Modern Roman}
\usepackage{amsmath,amssymb,mathtools}
\usepackage{unicode-math}
\setmathfont{Latin Modern Math}
\usepackage{xurl,hyperref}
\hypersetup{colorlinks=true,urlcolor=blue}
\setcounter{secnumdepth}{0}
\setlength{\parindent}{0pt}
\setlength{\parskip}{5pt}
\setlength{\emergencystretch}{3em}
\begin{document}
\section*{\texorpdfstring{Greenberg's Conjecture on Vanishing of Iwasawa Invariants of Totally Real Fields}{Greenberg's Conjecture on Vanishing of Iwasawa Invariants of Totally Real Fields}}\label{428}
\subsection{Definitions and mathematical statement}
Let $F$ be a totally real number field and $p$ a prime. Its cyclotomic ${\mathbb{Z}}_p$-extension is the infinite extension $F_\infty/F$ with Galois group ${\mathbb{Z}}_p$ obtained from the cyclotomic tower. Let $F_n$ be its degree-$p^n$ layer and $A_n$ the $p$-primary subgroup of the ideal class group of $F_n$. The Iwasawa invariants $\mu,\lambda$ are the nonnegative integers in the eventual formula
\[
 \log_p|A_n|=\mu p^n+\lambda n+\nu\qquad(n\gg0),
\]
where $\nu$ is an integer. Greenberg's selected conjecture is $\mu=\lambda=0$ for every such $F,p$, equivalently $\sup_n|A_n|<\infty$. It asserts more than the vanishing of $\mu$ alone, and the total-reality hypothesis is essential to this formulation.
\par\smallskip\textit{Formulation and concept references:} \href{https://arxiv.org/html/2503.15727v2}{[S1]}.

\subsection{Short English statement}
In a totally real field's cyclotomic tower, do the p-parts of the ideal class groups remain bounded in size?
\subsection{Research attempt}
\textit{GPT-6 Astra Ultra; 27 September 2026. Status: UNRESOLVED.}

This attempt investigates what finite class-group data would actually
need to certify. The target is bounded order, which requires both
$\mu=0$ and $\lambda=0$. I derive exact tests inside the eventual
range and then construct algebraic models separating the two growth
mechanisms. The models are not asserted to arise from totally real
number fields; their role is to expose what an arithmetic proof still
has to exclude.

\paragraph{Finite differences after the onset of the formula.}
Write $e_n=\log_p|A_n|$. If the Iwasawa formula holds at $n,n+1,n+2$,
then, with $\Delta e_n=e_{n+1}-e_n$,
\[
 \Delta e_n=\mu(p-1)p^n+\lambda,
\]
and hence
\[
 \mu=\frac{\Delta e_{n+1}-\Delta e_n}{(p-1)^2p^n},
 \qquad
 \lambda=\frac{p\Delta e_n-\Delta e_{n+1}}{p-1}.       \tag{1}
\]
These are exact integer identities. They provide consistency checks
for proposed eventual data, including divisibility and nonnegativity
conditions. In particular, a single flat step $e_{n+1}=e_n$ in a range
where the formula is known to hold forces $\mu=\lambda=0$, because
both summands in $\mu(p-1)p^n+\lambda$ are nonnegative. It then forces
all later orders to be constant.

The qualification about the range is indispensable. Given any cutoff
$N$, consider abstract finite groups
\[
 B_n=\mathbb Z/p^{\max(0,n-N)}\mathbb Z
\]
with the natural reduction maps $B_{n+1}\to B_n$. Their orders are
constant through level $N$, every transition map is surjective, and
their inverse limit is $\mathbb Z_p$. For $n\ge N$ their exponents
equal $n-N$, so the eventual invariants of their size sequence are
$\mu=0,\lambda=1,\nu=-N$. Arbitrarily long initial stabilization and
surjective transitions therefore do not imply eventual boundedness
for a general inverse system. This is a countertest to that inference,
not a construction of class groups of a number field.

\paragraph{Three explicit Iwasawa-algebra models.}
Let $\Lambda=\mathbb Z_p[[T]]$ and
$\omega_n=(1+T)^{p^n}-1$. For a $\Lambda$-module $M$, the quotient
$M/\omega_n M$ is its coinvariant module for the corresponding open
subgroup of the procyclic group. The following computations isolate
different possible sources of growth without invoking a classification
of arithmetic Iwasawa modules.

First take $M_\mu=\Lambda/(p)$. In characteristic $p$, repeated
Frobenius gives $(1+T)^{p^n}-1=T^{p^n}$. Thus
\[
 M_\mu/\omega_nM_\mu\cong
 \mathbb F_p[[T]]/(T^{p^n}),\qquad
 \log_p|M_\mu/\omega_nM_\mu|=p^n.                  \tag{2}
\]
This is the pure $\mu=1$ growth pattern. The module is torsion over
$\Lambda$, but torsion over this two-dimensional ring does not mean
finite as an abelian group.

For odd $p$, take $M_\lambda=\Lambda/(T-p)\cong\mathbb Z_p$.
Evaluation at $T=p$ is well-defined because the power series converges
$p$-adically. Its coinvariants have order
\[
 p^{v_p((1+p)^{p^n}-1)}=p^{n+1}.                    \tag{3}
\]
For completeness, if $v_p(u)\ge1$ and $p$ is odd, the binomial
expansion of $(1+u)^p-1$ has first term $pu$ of valuation
$v_p(u)+1$, while every other term has strictly larger valuation.
Induction starting at $u=p$ proves (3). For $p=2$, use
$\Lambda/(T-4)$ instead. If $v_2(u)\ge2$, then $2u+u^2$ has valuation
$v_2(u)+1$, giving order $2^{n+2}$. These models have $\mu=0$ but
$\lambda=1$; proving only the first invariant vanishes leaves a real
logical possibility of unbounded growth.

Finally let $M_0=\mathbb Z/p^a\mathbb Z$ with $T$ acting by zero.
Then $\omega_n$ acts by zero for every $n$, so all coinvariants have
order $p^a$. This is the bounded pattern. Direct sums multiply the
orders, hence add the exponents in (2)--(3). In particular the formal
sequence $p^n+n+1+a$ occurs from an explicit torsion module and has
both positive invariants. Formula (1) recovers them exactly.

\paragraph{What class-group ranks can and cannot show.}
The groups in (3) are cyclic: their $p$-ranks equal one at every level
while their orders tend to infinity. Consequently bounded
$\mathbb F_p$-dimension of $A_n/pA_n$ is insufficient for Greenberg's
claim. Likewise identifying an inverse limit as a finitely generated
$\mathbb Z_p$-module is not enough; a nonzero free $\mathbb Z_p$
summand is infinite and is compatible with the second model.
An argument using generators needs an independent uniform bound on
their orders, not merely a uniform bound on their number.

It would also be unjustified to identify $A_n$ with coinvariants of an
inverse-limit module without a control theorem. Norm maps, ramified
primes, and the discrepancy between a finite layer and its image in a
limit must be accounted for. The computations above are deliberately
stated for actual module quotients, so that no unproved arithmetic
control assertion is hidden in them.

\paragraph{A concrete certificate that would suffice.}
For a given pair $(F,p)$, one possible proof route would establish an
explicit index $n_0$ from which the Iwasawa formula applies and verify
one equality $|A_{n_0+1}|=|A_{n_0}|$. Equation (1) then finishes that
input. A different route could prove a norm-control statement with
uniformly bounded kernels and cokernels together with finiteness of
the relevant inverse limit. Neither route is supplied by a plateau
alone. The first gap here is a general arithmetic stabilization or
control theorem for every totally real $F$ and every prime $p$.

\paragraph{Source check, verification, and outcome.}
The recent paper [S1] treats specified real biquadratic families and
conditions, particularly at the prime two; its stated family results
do not remove the universal field and prime quantifiers. Exact finite
checks verify the valuations in (3), the characteristic-$p$ binomial
identity in (2), and the recovery formulas (1) on the explicit models.
No actual number-field class groups were computed in this attempt.
The outcome is a rigorous diagnostic for growth and stabilization,
including a demonstration that bounded rank and $\mu=0$ do not by
themselves answer the question. The stated Greenberg conjecture remains
\textbf{UNRESOLVED}.

\subsection{Sources}
\begin{itemize}
\item[S1]Mohamed Mahmoud Chems-Eddin et al., "Greenberg's conjecture and Iwasawa module of Real biquadratic fields I," arXiv:2503.15727 (2025); companion published version in J. Number Theory (2025).
\url{https://arxiv.org/html/2503.15727v2}.
\textit{Formulation source.}
\item[E1]Mohamed Mahmoud Chems-Eddin et al., \emph{Greenberg's conjecture
and Iwasawa module of Real biquadratic fields I}, arXiv:2503.15727v2,
Theorem 1.4 and its field and prime hypotheses.
\url{https://arxiv.org/html/2503.15727v2}.
\textit{Version-specific scope check of [S1]. Abstract module examples
here are not computations of the paper's number fields.
Consulted 27 September 2026.}

\end{itemize}
\end{document}
